\documentclass[10pt]{article}
\usepackage[paperwidth=8.5in,paperheight=19.6in,margin=0.45in]{geometry}
\usepackage{amsmath,amssymb,booktabs}
\pagestyle{empty}\setlength{\parindent}{0pt}\setlength{\parskip}{2pt}
\renewcommand{\arraystretch}{1.05}
\begin{document}
{\large\bfseries Current model (ordinary-matter shell + CoE core) vs.\ $A_F$, term by term}\hfill{\footnotesize worked sheet, 27 Sep 2026}\\[-2pt]
\rule{\textwidth}{0.4pt}
{\small
\textbf{1. Model and where each term comes from.} Sedov--Taylor flow $\vec V=k\,\vec x$ (position from the CoE; Taylor 1950; Sedov 1959), decelerated by a point core $M_c$ at the CoE and a uniform shell $\Omega_s$ between $R_{\rm in}$ and $R_{\rm out}$ (fixed in space; Earth and all sources inside it). Spherical symmetry makes the field radial, $\vec V=h(q)\,\vec q$, with $q=s/D$ ($D$ = CoE--Earth distance). Linear perturbation of the code's integration (from $\tau_i=0.01$, $H t\simeq1$):
\[
\frac{h(q)}{k}\simeq1-\tfrac12\Big[\underbrace{\tfrac{M_c}{M_{\rm crit}(D)}\big(\tfrac1{\tau_i}-1\big)q^{-3}}_{\text{core tidal term}}+\Omega_s\Big(\underbrace{\tfrac12}_{\text{uniform: }m=0\text{ only}}\ \underbrace{-\tfrac32 r_{\rm in}^2q^{-2}+r_{\rm in}^3q^{-3}}_{\text{hole }(r_{\rm in}=R_{\rm in}/D)}\Big)\Big],\qquad
\beta\equiv\frac{h'}{h}\Big|_{q=1},\ \ \gamma\equiv\frac{h''}{h}\Big|_{q=1}.
\]
Chris's $A$: $a_f=|\Delta\vec V|/{\rm separation}$ with sources at $R=14.3$~Gpc, $\varepsilon\equiv R/D$, $\mu=\cos b_0\cos\psi$, $\psi=\ell-\ell_0$ ($\ell_0$ = anti-CoE direction). Exactly $a_f\propto|h(q_s)(\hat x+\varepsilon\hat n)-h(1)\hat x|/\varepsilon$; to first order in $\varepsilon$ (the higher orders are the $r/D$ terms):
\[
a_f^2\propto1+B\mu^2+\varepsilon\,(P\mu+Q\mu^3),\qquad B=2\beta+\beta^2,\ \ P=\beta(3+\beta),\ \ Q=(1+\beta)(\gamma+\beta)-2\beta .
\]
Take the square root, then use $\cos^2\psi=\tfrac12(1+\cos2\psi)$, $\cos^3\psi=\tfrac14(3\cos\psi+\cos3\psi)$, $\cos^4\psi=\tfrac38+\tfrac12\cos2\psi+\tfrac18\cos4\psi$, $\cos^5\psi=\tfrac1{16}(10\cos\psi+5\cos3\psi+\cos5\psi)$ ($b_0=0$ as in the code):
\[
\frac{a_f}{\langle a_f\rangle}=1+a_1\cos\psi+a_2\cos2\psi+a_3\cos3\psi+\dots,\quad
a_2=\frac{\frac B4-\frac{B^2}{16}}{m_0},\ \
a_1=\frac{\varepsilon}{2m_0}\Big[P+\tfrac34\big(Q-\tfrac{PB}2\big)-\tfrac{5QB}{16}\Big],\ \
a_3=\frac{\varepsilon}{2m_0}\Big[\tfrac14\big(Q-\tfrac{PB}2\big)-\tfrac{5QB}{32}\Big],
\]
with $m_0=1+\frac B4-\frac{3B^2}{64}$. For a pure point mass ($\gamma=-4\beta$, $\beta\to3K$) this is the earlier line-of-sight form $K[(3\mu^2-1)+\varepsilon(\frac92\mu-\frac{15}2\mu^3)]$. The geometry gives only $\cos m\psi$ terms: even harmonics set by $\beta$, odd ones by $\varepsilon Q$. \emph{Rotation} (the sine factor, $A=a_f[1-{\rm Amp}\sin(\psi+\varphi)]$) adds
$C_1\!=\!a_1-{\rm Amp}\sin\varphi(1+\tfrac{a_2}2)$, $S_1\!=\!-{\rm Amp}\cos\varphi(1-\tfrac{a_2}2)$, $C_3\!=\!a_3-\tfrac{{\rm Amp}\,a_2}2\sin\varphi$, $S_3\!=\!-\tfrac{{\rm Amp}\,a_2}2\cos\varphi$ ($C_m,S_m$: $\cos m\psi,\sin m\psi$ parts).
\textbf{Latitude:} the code evaluates the model in the Galactic plane only (one direction per degree, no $|b|$ cut, no sector average), so the match is done the same way. A tilt $b_0$ would give $B\to B\cos^2b_0$ and $\varepsilon$-terms $\propto\cos b_0,\cos^3b_0$, degenerate with $\beta$ and $\varepsilon Q$; therefore $b_0=0$.

\textbf{2. $A_F$ rotated onto the model axis} ($\ell_0=14.96^\circ$, the $2\ell$ peak), $c_m\cos m\ell+s_m\sin m\ell=C_m\cos m\psi+S_m\sin m\psi$:
\begin{center}\begin{tabular}{lcccccccc}\toprule
 & $A_0$ & $A_1$, $\ell_1$ & $A_2$, $\ell_2$ & $A_3$, $\ell_3$ & $C_1$ & $S_1$ & $C_3$ & $S_3$\\\midrule
$A_F$ & 0.9965 & 0.0525, 55.7$^\circ$ & 0.0734, 14.96$^\circ$ & 0.0320, 67.0$^\circ$ & $0.0398$ & $0.0343$ & $-0.0292$ & $0.0130$\\\bottomrule
\end{tabular}\end{center}
\vspace{-4pt}
\textbf{3. Term-by-term solution} (Newtonian, series above, point-core shape $\gamma=-4\beta$; exact least-squares in brackets):
\begin{center}\footnotesize\setlength{\tabcolsep}{3pt}\begin{tabular}{llll}\toprule
Term & $A_F$ coefficient & Analytic expression & Solved value\\\midrule
$2\ell$ phase & $\ell_2=14.96^\circ$ & $\ell_0$ (axis; $\beta>0$: fast along it) & $\ell_0=14.96^\circ$ \ (12.8$^\circ$)\\
$2\ell$ amplitude & $C_2=0.0734$ & $a_2(\beta)$ & $\beta=0.161$ \ (0.210)\\
$1\ell$ (rotation) & $C_1=0.0398$, $S_1=0.0343$ & $C_1,S_1$ of Sec.~1 & Amp $=-0.068$, $\varphi=1.02$ \ (-0.072, 1.03)\\
$3\ell$ cos & $C_3=-0.0292$ & $a_3(\varepsilon Q)-\frac{{\rm Amp}\,a_2}2\sin\varphi$ & $\varepsilon Q=-0.308$ $\Rightarrow$ $\varepsilon=0.348$, $D=41$~Gpc \ (47)\\
$3\ell$ sin & $S_3=0.0130$ & $-\frac{{\rm Amp}\,a_2}2\cos\varphi$ (check) & predicted $0.0013$: \textbf{not produced}\\
$m=0$ & $A_0=0.9965$ & $N[1-\frac12{\rm Amp}\,a_1\sin\varphi]$ & $N=0.9971$\\\midrule
mass & -- & $\beta\simeq\frac{3K}{1-K}$, $K=\frac12\frac{M_c}{M_{\rm crit}(D)}(\frac1{\tau_i}-1)$ & $M_c/M_{\rm crit}(D)\approx0.0010$ (linear; $\approx$0.0017 nonlinear)\\\bottomrule
\end{tabular}\end{center}
\vspace{-4pt}
The $3\ell$ step fixes only $\varepsilon Q$: with the $0.909$-fit shape ($\gamma/\beta=-14.5$) it gives $D=114$~Gpc (exact 135), with the $\Omega_s=0.049$ shape ($\gamma/\beta=-3.7$) $D=39$~Gpc (exact 45). Every Newtonian solution reaches $R^2=0.987$ vs $A_F$ ($0.905$ vs Migkas).

\textbf{4. Cross-check: analytic series vs exact numeric projection of Chris's code} (geometric part, $a_m$ of $a_f/\langle a_f\rangle$; the code's own $\beta,\gamma$ from $h(q)$):
\begin{center}\footnotesize\setlength{\tabcolsep}{3pt}\begin{tabular}{lcccccccc}\toprule
Fit & $D$ [Gpc] & $\varepsilon$ & $\beta$ ($\beta_{\rm lin}$) & $\gamma$ & $a_1$ series / exact & $a_2$ series / exact & $a_3$ series / exact & $a_4$ exact\\\midrule
Newt.\ $\Omega_s$ free & 93.7 & 0.153 & 0.181 (0.26) & $-2.626$ & $-0.1082$ / $-0.0844$ & 0.0813 / 0.0807 & $-0.0448$ / $-0.0347$ & 0.0027\\
Newt.\ $\Omega_s$=0.049 & 49.2 & 0.291 & 0.157 (0.26) & $-0.575$ & $-0.0110$ / $-0.0126$ & 0.0717 / 0.0764 & $-0.0240$ / $-0.0256$ & 0.0069\\
Relativistic & 32.9 & 0.036$^a$ & 0.186$^b$ & -- & 0 / $-0.0043$ & 0.0849 / 0.0768 & 0 / $-0.0176$ & 0.0021\\\bottomrule
\end{tabular}\end{center}
{\footnotesize $^a$ $r/D$, $r=1.19$ Gpc. $^b$ $H_r/H_\perp-1$ at Earth.}\par\vspace{-2pt}
The linear formula for $h(q)$ overestimates $\beta$ by about 1.5--1.7 (the core pull at $\tau\sim\tau_i$ is not small), so $\beta,\gamma$ are read from the code; with them the series agrees with the exact projection to $\lesssim$30\% ($a_2$ to 6\%). Using the exact radial field $h(q)$ reproduces the code to $10^{-4}$. For the relativistic fit $A\propto D_L^{-1}$ at $z=0.3$ (sampling $r=1.19$~Gpc); the local shear gives $a_2\simeq\frac{\beta/2}{1+\beta/2}$ with $\beta=H_r/H_\perp-1$, and $r/D=0.036$ would make $a_1,a_3\approx0$, but the exact LTB curve has $a_3=-0.018$ (light-cone/shell-edge structure, not $r/D$).

\textbf{5. The three numerical fits, harmonic by harmonic} (full model with sine; amplitude @ peak longitude):
\begin{center}\footnotesize\setlength{\tabcolsep}{4pt}\begin{tabular}{lllllll}\toprule
 & $R^2$ Migkas & $R^2$ vs $A_F$ & $A_1$ & $A_2$ & $A_3$ & $M_c/M_{\rm crit}(D)$, $D$\\\midrule
$A_F$ & -- & 1 & 0.0525 @ 56$^\circ$ & 0.0734 @ 15$^\circ$ & 0.0320 @ 67$^\circ$ & --\\
Newtonian, $\Omega_s$ free & 0.909 & 0.987 & 0.0527 @ 56$^\circ$ & 0.0737 @ 12$^\circ$ & 0.0298 @ 72$^\circ$ & 0.0027, 93.7 Gpc\\
Newtonian, $\Omega_s$=0.049 & 0.893 & 0.977 & 0.0530 @ 55$^\circ$ & 0.0754 @ 13$^\circ$ & 0.0236 @ 72$^\circ$ & 0.0017, 49.2 Gpc\\
Relativistic, all free & 0.885 & 0.963 & 0.0528 @ 56$^\circ$ & 0.0763 @ 14$^\circ$ & 0.0160 @ 72$^\circ$ & 0.083, 32.9 Gpc\\\bottomrule
\end{tabular}\end{center}
\vspace{-4pt}
All three get $A_1$ (via the rotation) and $A_2$ (via $\beta$) right; they differ only in $A_3$, which ranks them (0.909 $>$ 0.893 $>$ 0.885). The hand solution ($\beta\approx0.16$, Amp $\approx-0.07$, $\varphi\approx1.0$, $\ell_0\approx13$--15$^\circ$, $\varepsilon Q\approx-0.31$) matches the $\Omega_s$=0.049 fit ($\beta=0.157$, Amp $=-0.062$, $\varphi=0.93$, $\ell_0=13^\circ$, $D=49$ vs 41--45~Gpc).

\textbf{6. Degeneracies and what the geometry cannot do.}
(i) \emph{Only $\beta$ is fixed by $2\ell$.} $\beta$ depends on $M_c/M_{\rm crit}(D)$, $\Omega_s$ and $R_{\rm in}/D$, all ratios: the core's tidal field in units of $H^2$ is $GM_c/(H^2D^3)=\frac12M_c/M_{\rm crit}(D)$. $D$ enters only through $\varepsilon=r/D$ in the odd terms.
(ii) \emph{Newtonian:} $r=R=14.3$~Gpc, so $3\ell$ fixes $\varepsilon Q$. $D$ trades against the curvature $\gamma$ of $h(q)$, i.e. the shell structure ($D\approx41$--47, 39--45 and 114--135~Gpc for the point-core, $\Omega_s$=0.049 and 0.909 shapes, same $R^2$).
(iii) \emph{Relativistic:} $r=1.19$~Gpc, so $r/D\le0.036$ and the odd terms do not depend on $D$. Along the fitted valley $a_2$ = 0.076--0.078 and $a_3$ = $-$0.016--$-$0.018 stay the same for $D$ = 33--$7\times10^6$~Gpc while $M_c/M_{\rm crit}(D)$ = 0.061--0.083. Hence the LTB fit pins $M_c/M_{\rm crit}(D)$, not $D$.
(iv) \emph{Not produced by the geometry:} the dipole is along the axis ($a_1\cos\psi$, peak at $\ell_0$ or $\ell_0+180^\circ$), but $A_F$ peaks at $\ell_1=55.7^\circ$, so the off-axis dipole needs the rotation (sine). The $3\ell$ phase offset ($S_3=+0.013$ vs $\le0.002$ from rotation) cannot be made: the model's $3\ell$ peak is locked to $\ell_0+60^\circ$ (data $\ell_3=67^\circ$ vs $\ell_0+60\approx73$--$75^\circ$). This missing ``$3\ell$ transform'' is the main residual of every fit.


\textbf{7. Input parameters implied by the term-by-term solution vs.\ Table~1} of universe\_bestfit\_thinshell.pdf. Each column keeps that fit's shell assumption. Newtonian $\Omega_s$ free: $\Omega_s=1.012$, thin shell $R_{\rm in}=D-R$, $R_{\rm out}=D+R$. Newtonian $\Omega_s$=0.049: $R_{\rm in}=D-1.028\,R$, $R_{\rm out}=D+1.475\,R$. Relativistic: $R_{\rm in}=D-1.040\,R$, $R_{\rm out}=D+1.171\,R$. The unknowns are solved from the Sec.~3 equations ($2\ell$ phase $\to\ell_0$, $2\ell$ amplitude $\to\beta\to M_c/M_{\rm crit}$, $1\ell\to$ Amp, $\varphi$, $3\ell$ cos $\to\varepsilon Q\to D$; $C_m,S_m$, $m=1$--3, by least squares). The $a_m(\beta,\gamma)$ come from the code's exact $h(q)$, not the linear formula. Masses use $H_0=70$, $\rho_s=\Omega_s\rho_c$.
\begin{center}\footnotesize\setlength{\tabcolsep}{2.6pt}\begin{tabular}{lllllll}\toprule
 & \multicolumn{2}{c}{Newtonian, $\Omega_s$ free} & \multicolumn{2}{c}{Newtonian, $\Omega_s$=0.049} & \multicolumn{2}{c}{Relativistic, all free}\\
Quantity & term-by-term & Table 1 & term-by-term & Table 1 & term-by-term (Newt., $r$=1.19) & Table 1 (LTB)\\\midrule
$D/R$ & 6.474 & 6.554 & 2.589 & 3.439 & \multicolumn{1}{c}{not fixed} & 2.303\\
$D$ & $2.857{\times}10^{24}$ km & $2.892{\times}10^{24}$ km & $1.143{\times}10^{24}$ km & $1.517{\times}10^{24}$ km & any $D$ (tried 33--330) & $1.016{\times}10^{24}$ km\\
 & (92.6 Gpc) & (93.7 Gpc) & (37.0 Gpc) & (49.2 Gpc) &  & (32.9 Gpc)\\
$R_{\rm in}$ & 78.3 Gpc & 79.4 Gpc & 22.3 Gpc & 34.5 Gpc & $D-14.9$ Gpc & 18.1 Gpc\\
$R_{\rm out}$ & 106.9 Gpc & 108.0 Gpc & 58.1 Gpc & 70.3 Gpc & $D+16.7$ Gpc & 49.7 Gpc\\
$M_c$ & $2.38{\times}10^{54}$ kg & $2.47{\times}10^{54}$ kg & $9.41{\times}10^{52}$ kg & $2.31{\times}10^{53}$ kg & $\propto D^3$ ($7.1{\times}10^{52}$ kg at 32.9 Gpc) & $3.38{\times}10^{54}$ kg\\
 & $1.20{\times}10^{24}$ $M_\odot$ & $1.24{\times}10^{24}$ $M_\odot$ & $4.73{\times}10^{22}$ $M_\odot$ & $1.16{\times}10^{23}$ $M_\odot$ &  & $1.70{\times}10^{24}$ $M_\odot$\\
$M_c/M_{\rm crit}(D)$ & 0.00265 & 0.00265 & 0.00164 & 0.00171 & \textbf{0.0016--0.0018}$^a$ & 0.0835\\
$M(<D)$, total & $3.62{\times}10^{56}$ kg & $3.72{\times}10^{56}$ kg & $2.29{\times}10^{54}$ kg & $4.56{\times}10^{54}$ kg & $\propto D^3$ & $5.04{\times}10^{54}$ kg\\
$M_{\rm shell}$, total & $8.50{\times}10^{56}$ kg & $8.71{\times}10^{56}$ kg & $1.03{\times}10^{55}$ kg & $1.70{\times}10^{55}$ kg & $\propto D^2$ & $6.48{\times}10^{54}$ kg\\
$M$ within $R$ & $3.35{\times}10^{54}$ kg & $3.35{\times}10^{54}$ kg & $1.62{\times}10^{53}$ kg & $1.62{\times}10^{53}$ kg & $1.62{\times}10^{53}$ kg & $1.62{\times}10^{53}$ kg\\
$\Omega_s$ & 1.012$^b$ & 1.012 & 0.049$^b$ & 0.049 & 0.049$^b$ & 0.049\\
$\ell_0$ & 13.2$^\circ$ & 12.8$^\circ$ & 12.8$^\circ$ & 13.0$^\circ$ & 15.0$^\circ$ & 13.5$^\circ$\\
Amp & $-0.1246$ & $-0.1240$ & $-0.0637$ & $-0.0622$ & $-0.0528$ & $-0.0558$\\
$\varphi$ & 1.267 rad & 1.264 rad & 0.941 rad & 0.929 rad & 0.819 rad & 0.848 rad\\
$R^2$ vs Migkas & 0.909 & 0.909 & 0.881 & 0.893 & 0.824$^c$ & 0.885\\
$R^2$ vs $A_F$ & 0.987 & 0.987 & 0.970 & 0.977 & 0.899$^c$ & 0.963\\
\bottomrule
\end{tabular}\end{center}
\vspace{-3pt}
{\footnotesize $^a$ Newtonian code, $\tau_i=0.01$, sampling $r=1.19$ Gpc: the $r/D$-odd terms vanish, so $D$ is not fixed; any $D$ works with $M_c/M_{\rm crit}(D)\approx0.0017$ (0.0018, 0.0017, 0.0016 at $D$ = 33, 100, 330 Gpc). The exact LTB curve instead gives $M_c/M_{\rm crit}(D)=0.061$--0.083 for any $D$ (Sec.~6). $^b$ Held fixed (that column's shell assumption). $^c$ At $D=32.9$ Gpc. The harmonic solve weights $C_m,S_m$ equally, not the 71 points, so its $R^2$ can sit slightly below the point fit.}

The \textbf{0.909 fit} is recovered: $D$ 92.6 vs 93.7 Gpc, the same $M_c/M_{\rm crit}$, and $\ell_0$, Amp, $\varphi$ to 3\%. For \textbf{$\Omega_s$=0.049} the harmonic solve gives the same $M_c/M_{\rm crit}$ (0.00164 vs 0.00171), $\ell_0$, Amp and $\varphi$, but a smaller $D$ (37 vs 49 Gpc): the harmonic weighting pulls toward a stronger $3\ell$ term.
\textbf{$\tau_i$ test} (Chris's hypothesis that the code pins $M_c/M_{\rm crit}\cdot(1/\tau_i-1)$). The $\Omega_s$=0.049 fit was redone at other start times, with $D$, $M_c$, both shell edges, $\ell_0$, Amp and $\varphi$ free:
\begin{center}\footnotesize\begin{tabular}{lcccc}\toprule
$\tau_i$ & $M_c/M_{\rm crit}(D)$ & $M_c/M_{\rm crit}\cdot(1/\tau_i-1)$ & $D$ [Gpc] & $R^2$\\\midrule
0.1 & 0.0112 & 0.101 & 55.4 & 0.894\\
0.01 (Table 1) & 0.0017 & 0.169 & 48.7 & 0.893\\
0.001 & 0.0015 & 1.491 & 47.0 & 0.892\\\bottomrule
\end{tabular}\end{center}
\vspace{-3pt}
The hypothesis holds only partly. From $\tau_i$ = 0.01 to 0.1, $M_c/M_{\rm crit}$ rises 6.5$\times$ (11$\times$ predicted). From 0.01 to 0.001 it falls only 13\% (10$\times$ predicted), because the core's own size ($L_c=0.005\,D$) cuts off the early pull. So the Newtonian $M_c/M_{\rm crit}$ is set by the code's start time and core size, and $M_c/M_{\rm crit}\cdot(1/\tau_i-1)$ is not constant (0.10--1.5). The $\sim$50$\times$ gap to the relativistic 0.083 is mostly this regularisation: at $\tau_i=0.1$ it shrinks to 7$\times$. The LTB value (single bang time, no $\tau_i$) is the physically calibrated one.

{\footnotesize\textbf{References.} Taylor, G.~I. 1950, Proc.\ R.\ Soc.\ Lond.\ A, 201, 159. Sedov, L.~I. 1959, Similarity and Dimensional Methods in Mechanics (Academic Press). Migkas, K., et al. 2020, A\&A, 636, A15. Scripts: A\_equals\_fourier/thinshell/tm\_core.py, termmatch.py, relvalley.py, make\_tex.py, solve\_inputs.py, taui\_test.py; numbers in termmatch.json, solve\_inputs.json, taui\_test.json.}
}
\end{document}
